% https://tex.stackexchange.com/questions/23647/drawing-a-directory-listing-a-la-the-tree-command-in-tikz
%
\newcommand{\FTdir}{}
\def\FTdir(#1,#2,#3){%
    \FTfile(#1,{{\color{black!40!white}\faFolderOpen}\hspace{0.2em}#3})
    (tmp.west)++(0.8em,-0.4em)node(#2){}
    (tmp.west)++(1.5em,0)
    ++(0,-1.3em) 
}

\newcommand{\FTfile}{}
\def\FTfile(#1,#2){%
    node(tmp){}
    (#1|-tmp)++(0.6em,0)
    node(tmp)[anchor=west,black]{\ttfamily #2}
    (#1)|-(tmp.west)
    ++(0,-1.2em) 
}

\newcommand{\FTroot}{}
\def\FTroot{tmp.west}

\subsection{Top-level structure}
\begin{tikzpicture}%
\draw[color=black!60!white]
\FTdir(\FTroot,root,GCWizard){       % root: parent = \FTroot
    
    \FTdir(root,android,android){       % normal dir: (parentID, currentID, label)
    }
    
%    ++(0,-0.5em)                % additional space if neded
    
    \FTdir(root,assets,assets) {
    }

%++(0,-0.5em)

    \FTdir(root,ios,ios) {
    }
    
%    ++(0,-0.5em)
    
    \FTdir(root,lib,lib){
    }

    \FTdir(root,test,test){
    }
};
\end{tikzpicture}
\captionof{figure}{Top-level structure}
\label{fig: structure-toplevel}

\begin{description}
    \item[\textbackslash android] contains the data for building the android app.
    \item[\textbackslash assets] contains the ressorce files e.g. icons, symbols, language-files.
    \item[\textbackslash ios] contains the data for building the iOS app.
    \item[\textbackslash lib] contains the source for the frontend and the backend.
    \item[\textbackslash test] contains the tests.
\end{description}

\subsection{assets-directory}
\begin{tikzpicture}%
\draw[color=black!60!white]
\FTdir(\FTroot,root,GCWizard){       % root: parent = \FTroot
    
    \FTdir(root,assets,assets) {
        \FTdir(assets,coordinates,coordinates){}
        \FTdir(assets,graphics,graphics){}
        \FTdir(assets,i18n,i18n){
            \FTfile(i18n,de.json)
            \FTfile(i18n,en.json)
        }
        \FTdir(assets,logo,logo){}
        \FTdir(assets,symboltables,symbol\_tables){
            \FTdir(symboltables,various,\dots)
        }
    }
};
\end{tikzpicture}
\captionof{figure}{Sub-directory assets}
\label{fig: structure-assets}

\begin{description}
    \item[\textbackslash coordinates] contains the icon-files for the coordinates menu.
    \item[\textbackslash graphics] contains the ressource files e.g. icons, symbols, language-files.
    \item[\textbackslash i18n] contains the language-files; currently English and German are supported. For other languages thes files have to be translated.
    \item[\textbackslash logo] contains the files for the GCWizard-logo.
    \item[\textbackslash smybol\_tables] contains all the symbols. For every symbol table exists a separate subdirectory.
\end{description}

\subsection{lib-directory}
\begin{tikzpicture}%
\draw[color=black!60!white]
\FTdir(\FTroot,root,GCWizard){       % root: parent = \FTroot
    
    \FTdir(root,lib,lib){
        \FTdir(lib,database,database){
            \FTdir(database,common,common)
            \FTdir(database,formulasolver,formula\_solver)
            \FTdir(database,logs,logs)
        }
        \FTdir(lib,i18n,i18n){}
        \FTdir(lib,logic,logic){
            \FTdir(logic,tools,tools){
                \FTdir(tools,coords,coords)
                \FTdir(tools,cryptoandencodings,crypto\_and\_encodings)
                \FTdir(tools,formulasolver,formula\_solver)
                \FTdir(tools,games,games)
                \FTdir(tools,scienceandtechnology,science\_and\_technology)
            }
            \FTdir(logic,units,units)
        }
        \FTdir(lib,persistence,persistence)
        \FTdir(lib,theme,theme)
        \FTdir(lib,utils,utils)
        \FTdir(lib,widgets,widgets){
            \FTdir(widgets,common,common){}
            \FTdir(widgets,mainmenu,main\_menu){}
            \FTdir(widgets,selectorlists,selector\_lists){}
            \FTdir(widgets,tools,tools){
                \FTdir(tools,coords,coords)
                    \FTdir(tools,cryptoandencodings,crypto\_and\_encodings)
                    \FTdir(tools,formulasolver,formula\_solver)
                    \FTdir(tools,games,games)
                    \FTdir(tools,scienceandtechnology,science\_and\_technology)
            }
            \FTdir(widgets,utils,utils){}
            \FTfile(widgets,favorites.dart)
            \FTfile(widgets,main\_menu.dart)
            \FTfile(widgets,main\_view.dart)
            \FTfile(widgets,registry.dart)
        }
        \FTfile(lib,main.dart)
    }
};
\end{tikzpicture}
\captionof{figure}{Sub-directory lib}
\label{fig: structure-lib}

\begin{description}
    \item[\textbackslash database] contains the source code for managing the internal database.
    \item[\textbackslash i18n] contains the source code for providing the internationalization.
    \item[\textbackslash logic] contains the source code for the backend.
    \item[\textbackslash persistence] contains the additional source code for managing formulas and calculating coordinates..
    \item[\textbackslash theme] contains the source code for design of the front end.
    \item[\textbackslash utils] contains the source code for common usable definitions.
    \item[\textbackslash widgets] contains the source code for creating the frontend.
\end{description}

\subsection{test-directory}
\begin{tikzpicture}%
\draw[color=black!60!white]
\FTdir(\FTroot,root,GCWizard){       % root: parent = \FTroot
    
    \FTdir(root,test,test){
        \FTdir(test,logic,logic) {
            \FTdir(logic,tools, tools){
                \FTdir(tools,coords,coords){}
                \FTdir(tools,cryptoandencodings,crypto\_and\_encodings)
                \FTdir(tools,formulasolver,formula\_solver)
                \FTdir(tools,scienceandtechnology,science\_and\_technology)
            }
        }
        \FTdir(test,utils,utils){}
    }
};
\end{tikzpicture}
\captionof{figure}{Sub-directory tests}
\label{fig: structure-tests}

\begin{description}
    \item[\textbackslash android] contains the data for building the android app.
    \item[\textbackslash assets] contains the ressource files e.g. icons, symbols, language-files.
    \item[\textbackslash ios] contains the data for building the iOS app.
    \item[\textbackslash lib] contains the source for the frontend and the backend.
    \item[\textbackslash test] contains the tests.
\end{description}

\subsection{Important directories and files for contributers}
\subsubsection{Directories}
\begin{description}
    \item[\textbackslash lib\textbackslash logic\textbackslash] This directory with its subdirectories provides all the stuff for the backend.
    \begin{description}
        \item[tools] This directory with its subdirectories provides all the stuff for the
        \begin{description}
            \item[coords] Methods algorithms to calculate coordinates
            \item[crypto\_and\_encodings] Methods for encryption/decryption or encoding/decoding
            \item[formula\_solver] Methods for the formula solver
            \item[games] Methods for encryption/decryption or encoding/decoding due to games like Scrabble, Sudoku etc.
            \item[science\_and\_technology] Methods and calculations for scientific stuff
        \end{description}
        \item[units] These files provide all units which can be used for calculation like in Converting Units, Windchill, Heat Index, Beaufort
    \end{description}
    \item[\textbackslash lib\textbackslash utils\textbackslash] This directory contains common used definitions like alphabets, constants etc.
    \item[\textbackslash lib\textbackslash widget\textbackslash] This directory with its subdirectories provides all the stuff for the frontend.
    \begin{description}
        \item[common] This directory with its subdirectories provides all the widgets for input and output.
        \item[selector\_lists] This directory contains all dart-files to create submenues.
        \item[tools] This directory with its subdirectories contains all the widgets.
        \begin{description}
            \item[coords] Widgets to calculate coordinates
            \item[crypto\_and\_encodings] Widgets for encryption/decryption or encoding/decoding
            \item[formula\_solver] Widgets for the formula solver
            \item[games] Widgets for encryption/decryption or encoding/decoding due to games like Scrabble, Sudoku etc.
            \item[science\_and\_technology] Widgets and calculations for scientific stuff
        \end{description}
    \end{description}
\end{description}

\subsubsection{Files}
\begin{description}
    \item[\textbackslash assets\textbackslash i18n\textbackslash de.json] Internationalization for German language.
    \item[\textbackslash assets\textbackslash i18n\textbackslash en.json] Internationalization for English language.
    \item[\textbackslash lib\textbackslash widgets\textbackslash main\_view.dart] Register all classes
    \item[\textbackslash lib\textbackslash widgets\textbackslash registry.dart] Define all classes with their metadata
    \begin{description}
        \item[tool] Name of the Class.
        \item[i18nPrefix] links the tool to the internationalization.
        \item[category] ToolCategory.CRYPTOGRAPHY, .SCIENCE\_AND\_TECHNOLOGY,
        \item[searchStrings] Strings for the search - with these keywords the tool can be found.
        \item[iconPath] if there should be an icon displayed, the paths has to be entered here.
    \end{description}
\end{description}


